\scanpage{34}

\deleted{entiers $i,j,l$, donnés par leurs développements $p$-adiques}

Soit $r\geq 0$ tel que $p^r>i+j$, soit \deleted{\unsure{$l=$}}

\[
l=p^r-1-(i+j)
\]

\deleted{Je dis que} de sorte que l'on a

\[
\underset{\substack{\rotatebox{25}{$=$}\\\inserted{(i+j)+l=i+(j+l)}}}{i+j+l}=p^r-1=\sum_{\alpha=0}^{r-1}(p-1)p^\alpha
\]

On sait alors, par le critère habituel, que l'on a $C_{i+j,l}\not\equiv 0\ (p)$, et pour la $=$ raison on a $C_{\unsure{j|i},j+l}\not\equiv 0\ (p)$

cqfd

\uline{NB} Le lemme commun qui a servi dans les lemmes 1 et 2 est le suivant :

Si $\mu,\nu\geq 0$ sont tels que

\[
\mu+\nu=p^r-1\qquad (r\geq 0\ \text{convenable})
\]

\[
\text{alors}\qquad C_{\mu,\nu}\not\equiv 0\ (p)
\]

\deleted{(plus précisément, $C_{\mu,\nu}$}

% archive sheet number: 33
